% Vietnamese translation for OI Public Library
% Source-File: IOI中国国家候选队论文/2006/朱泽园/朱泽园.ppt
\documentclass[11pt]{article}
\usepackage{oipl-vn}

\newcommand{\Oh}{\mathcal{O}}

\title{Half-plane Intersection\\{\large Ghi chú theo slide}}
\author{Zeyuan Zhu}
\date{2006}

\begin{document}
\maketitle

\origtitle{Half-plane Intersection and practical algorithms}
\sourcepath{IOI China candidate team papers / 2006 / Zhu Zeyuan / Zhu Zeyuan.ppt}

\section{Phạm vi}

Bộ trình chiếu đi kèm bài viết về giao nửa mặt phẳng. Ghi chú này dựa trên
bản DOC/PDF đã dịch, tập trung vào các khái niệm và thuật toán mà slide dùng
để dẫn dắt.

\section{Bài toán}

Một nửa mặt phẳng có dạng
\[
  ax+by\le c.
\]
Cho nhiều nửa mặt phẳng, cần tìm miền điểm thỏa tất cả bất đẳng thức. Miền
kết quả là lồi, có thể rỗng, suy biến, bị chặn hoặc không bị chặn. Trong cài
đặt thi đấu, người ta thường thêm một hình chữ nhật rất lớn để làm miền xét
bị chặn.

\section{Giao đa giác lồi}

Bài viết dùng giao hai đa giác lồi như bước chuẩn bị. Nếu ta biết duy trì
biên của miền khả thi sau mỗi lần thêm ràng buộc, giao với một nửa mặt phẳng
mới chỉ cần giữ các đoạn biên còn nằm trong phía hợp lệ và thêm các giao điểm
cắt biên.

\section{Hai hướng thuật toán}

Hướng chia để trị tách tập nửa mặt phẳng thành hai phần, giải từng phần rồi
giao hai đa giác lồi thu được. Hướng sắp xếp rồi thêm dần xử lý các nửa mặt
phẳng theo thứ tự góc, dùng hàng đợi hai đầu để loại các cạnh không còn nằm
trên biên kết quả. Trong trường hợp đã có thứ tự thuận lợi, phần thêm dần có
thể chạy tuyến tính.

\section{Giá trị thực hành}

Điểm mạnh của thuật toán là giữ dữ liệu hình học dưới dạng biên lồi rõ ràng.
Các tình huống song song, trùng biên và miền rỗng cần được xử lý nhất quán,
vì sai số số thực trong những trường hợp này thường gây lỗi khó phát hiện.

\end{document}
